\chapter{实证分析}
\label{chap:empirical}

\section{数据来源与说明}

本文使用虚构的 2018—2022 年省际面板数据进行演示，所有数值均为模拟生成，不代表真实情况。为便于比较，将样本地区划分为东部、中部与西部三组。

\section{综合测度结果}

按第~\ref{chap:method}~章方法测算，部分地区数字经济发展水平综合得分如表~\ref{tab:score} 所示。可以看到，东部地区得分整体高于中西部，且呈逐年上升趋势。

\begin{table}[!htbp]
    \centering
    \caption{部分地区数字经济发展水平综合得分（示例）}
    \label{tab:score}
    \wuhao\songti
    \begin{tabular}{lccc}
        \toprule
        地区 & 2018 年 & 2020 年 & 2022 年 \\
        \midrule
        A 省 & 0.412 & 0.498 & 0.576 \\
        B 省 & 0.305 & 0.361 & 0.423 \\
        C 省 & 0.241 & 0.287 & 0.339 \\
        \bottomrule
    \end{tabular}

    \dufesource{资料来源：本文基于虚构数据测算整理。}
\end{table}

\section{区域差异分解}

为识别区域差异的来源，本文先测算总体泰尔指数，再将其分解为组间与组内两部分。

\subsection{泰尔指数测算}

按式~\eqref{eq:theil} 计算总体泰尔指数，结果见表~\ref{tab:theil}。

\begin{table}[!htbp]
    \centering
    \caption{泰尔指数分解结果（示例）}
    \label{tab:theil}
    \wuhao\songti
    \begin{tabular}{clcc}
        \toprule
        年份 & 差异来源 & 贡献值 & 贡献率 \\
        \midrule
        \multirow{2}{*}{2018} & 组间差异 & 0.061 & 64\% \\
                              & 组内差异 & 0.034 & 36\% \\
        \midrule
        \multirow{2}{*}{2022} & 组间差异 & 0.078 & 69\% \\
                              & 组内差异 & 0.035 & 31\% \\
        \bottomrule
    \end{tabular}

    \dufesource{资料来源：本文基于虚构数据测算整理。}
\end{table}

\subsection{组间与组内分解}

将总差异分解后可见，区域差异主要来源于组间差距。综合得分的区域分布如图~\ref{fig:dist} 所示，直观印证了上述格局。

\begin{figure}[!htbp]
    \centering
    \includegraphics[width=0.8\textwidth]{figures/demo.png}
    \caption{综合得分的区域分布（示例）}
    \label{fig:dist}
\end{figure}
